% 引入必要的包
\usepackage{amsmath, amssymb, amsfonts,amsthm}
\usepackage{graphicx}
\usepackage{geometry}
\usepackage{fancyhdr}
\usepackage{titlesec}
\usepackage{hyperref}
\usepackage{xcolor}
\usepackage{listings}
\usepackage{booktabs}
\usepackage{caption}
\usepackage{enumitem}
\usepackage{tcolorbox}
\usepackage{eso-pic}
\usepackage{fontawesome5}
\usepackage{tocloft}
\usepackage{multirow}
\usepackage{float}

% 文献引用
\usepackage[utf8]{inputenc}
\usepackage[backend=biber,style=numeric,citestyle=numeric]{biblatex}  
\addbibresource{references.bib}  % 引入 references.bib 文件
\hypersetup{colorlinks=true, linkcolor=black, urlcolor=black, citecolor=black}

% 设置图片目录
\graphicspath{{Figures/}}

\definecolor{blackk}{HTML}{0d0d0d}
\definecolor{cyan}{rgb}{0.0, 1.0, 1.0}          % 青色
\definecolor{lightblue}{rgb}{0.68, 0.85, 0.9}   % 浅蓝色
\definecolor{darkgreen}{rgb}{0.0, 0.5, 0.0}     % 墨绿色
\definecolor{brownyellow}{rgb}{0.76, 0.60, 0.42} % 棕黄色

% 页面布局设置
\geometry{left=2.5cm, right=2.5cm, top=2.5cm, bottom=2.5cm}
\setlength{\headheight}{14.5pt}
\addtolength{\topmargin}{-2.5pt}

% 添加封面图片
% \AddToShipoutPicture*{
%     \put(0,0){%
%         \includegraphics[width=\paperwidth,height=\paperheight]{Figures/cover.jpg}}
% }

% 自定义标题样式
\titleformat{\section}
  {\normalfont\Large\bfseries\color{black}}    % 一级标题青色
  {\thesection}{1em}{}

\titleformat{\subsection}
  {\normalfont\large\bfseries\color{cyan}} % 二级标题浅蓝色
  {\thesubsection}{1em}{}

\titleformat{\subsubsection}
  {\normalfont\normalsize\bfseries\color{lightblue}} % 三级标题墨绿色
  {\thesubsubsection}{1em}{}

\titleformat{\paragraph}
  {\normalfont\normalsize\bfseries\color{darkgreen}} % 四级标题棕黄色
  {\theparagraph}{1em}{}

% 修改目录的字体颜色为黑色
\hypersetup{
    colorlinks=true,
    linkcolor=black,  % 将目录超链接设置为黑色
    urlcolor=magenta,
    citecolor=black,
    pdfborder={0 0 0}
}

% 页眉页脚设置
\pagestyle{fancy}
\fancyhf{}
\fancyhead[L]{\leftmark}
\fancyhead[R]{\thepage}

% \fancyfoot[R]{\href{https://github.com/hanlife02}{\faGithub\ Ethan}}

% 代码高亮设置
\lstset{
    basicstyle=\ttfamily\footnotesize,
    keywordstyle=\color{blue},
    commentstyle=\color{gray},
    stringstyle=\color{red},
    frame=single,
    breaklines=true
}

% 设置自定义样式的定理框
\tcbuselibrary{theorems}
% 定义定理的样式
\newtcolorbox{dl}[1][]{colback=red!5!white,colframe=red!75!black,
    fonttitle=\bfseries,title=定理, #1}

% 定义定义的样式
\newtcolorbox{dy}[1][]{colback=green!5!white,colframe=green!75!black,
    fonttitle=\bfseries,title=定义, #1}

% 定义公式的样式
\newtcolorbox{gs}[1][]{colback=blue!5!white,colframe=blue!75!black,
    fonttitle=\bfseries,title=公式, #1}

% 定义旁注的样式，支持传入参数作为标题
\newtcolorbox{sidenote}[2][]{title=#2, #1}

